\documentclass[10pt,a4paper]{amsart}
\usepackage[margin=19mm]{geometry}
\usepackage{amsmath,amssymb,mathtools}
\usepackage[scr=rsfs,cal=euler]{mathalfa}
\usepackage{xfrac}
\usepackage[unicode,hidelinks,bookmarks=false]{hyperref}
\newcommand{\ie}{i.e.\ }
\newcommand{\cf}{cf.\ }
\newcommand{\resp}{respectively\ }
\newcommand{\Subsection}{\subsection}
\providecommand{\nobreakdash}{\nobreak\hbox{-}}
\setlength{\emergencystretch}{2em}
\multlinegap0pt
\usepackage{amsthm}
\theoremstyle{plain}
\newtheorem{theorem}{Theorem}[section]
\newtheorem{corollary}[theorem]{Corollary}
\newtheorem{main}{Main Theorem}
\newtheorem{lemma}[theorem]{Lemma}
\newtheorem{proposition}{Proposition}
\newtheorem{conjecture}{Conjecture}
\newtheorem{problem}{Problem}
\newtheorem{definition}[theorem]{Definition}
\newtheorem{remark}[theorem]{Remark}
\newtheorem{notation}{Notation}
\newcommand{\ep}{\varepsilon}
\newcommand{\eps}[1]{{#1}_{\varepsilon}}
\newtheorem{example}[theorem]{Example}
\title[Lyapunov spectra behavior for linear discrete-time systems under small perturbations]{Lyapunov spectra behavior for linear discrete-time systems under small perturbations: prescribed shifts of the Lyapunov exponents realized by an admissible multiplicative perturbation (Theorem 5.1)}
\author{Adam Czornik, Evgenii Makarov, Michal Niezabitowski, Svetlana Popova, and Vasilii Zaitsev}
\date{Source: arXiv:2106.14691v1 (CC BY 4.0)}
\begin{document}
\maketitle
\noindent\emph{Source and licence.} Lyapunov spectra behavior for linear discrete-time systems under small perturbations, arXiv:2106.14691v1, distributed under the Creative Commons Attribution 4.0 International licence, http://creativecommons.org/licenses/by/4.0/, as recorded in the arXiv OAI metadata for this exact version and shown on the abstract page https://arxiv.org/abs/2106.14691v1 (v1 of 28 June 2021; the work appeared in IEEE Access, DOI 10.1109/ACCESS.2022.3189566). Full rights notice: licenses/arxiv-2106.14691-CC-BY-4.0.txt.\par\smallskip
\noindent\textit{Editorial scope.} Counted unit: the paper\textquotesingle s basic result of section 5, source label \textbackslash{}label\{T6\} (source0.tex lines 1104--1121), printed as Theorem 5.1 of the version PDF (page 13), delivered with its complete original proof as the single primary proof body. This article uses no proof environment at all: all seven of its proofs are run-in \textbackslash{}emph\{Proof.\} blocks that follow their statements directly. The counted proof is therefore delivered as the exact run-in block that belongs to this theorem, from its \textbackslash{}emph\{Proof.\} marker at source0.tex line 1123 to the \textbackslash{}hfill $\square$ that closes it at source0.tex line 1354, the whole 232-line argument up to but not including the following \textbackslash{}section\{Applications\}; it is one primary proof body and it contains the complete argument, with the numbered displays X2, X3, N7 and X6 and the unnumbered displays X4 and X5.\par
\noindent\textit{Required context retained uncounted.} Four source spans: the definition of the discrete-time system (4) (source0.tex lines 169--171, source label \textbackslash{}label\{2\}); the multiplicatively perturbed system (8) with the sentence defining a multiplicative perturbation (lines 519--524, source label \textbackslash{}label\{M4\}); Lemma 2.2 (lines 308--318, source label \textbackslash{}label\{L3\}); and Lemma 2.3 (lines 320--349, source label \textbackslash{}label\{L2\}). These are exactly the four objects the counted statement and proof name from outside their own text, and all four are needed for the argument to read as the author wrote it. No other source span is delivered. The author\textquotesingle s own remark, definitions and earlier sections are not reproduced; the notions of a Lyapunov sequence, of an FSS, of a splitted, broken-away, normal or equidistantly-angular family and of the Lyapunov exponent lambda[.] are used as the authors define them in the article\textquotesingle s sections 2--4, which lie outside this unit.\par
\noindent\textit{Wrapper and numbering.} The article\textquotesingle s \textbackslash{}documentclass[11pt]\{article\} and its package list are replaced by the AMS amsart wrapper of the pinned TeX Live runtime; the author\textquotesingle s own theorem environments (source lines 25--39) and his two macros \textbackslash{}ep and \textbackslash{}eps (lines 35--36) are carried over verbatim. The author\textquotesingle s own redefinitions of \textbackslash{}@begintheorem and \textbackslash{}@opargbegintheorem (source lines 18--22, a bold run-in head style) are not carried over, so the statements print in the wrapper\textquotesingle s standard AMS theorem style. Editorial counter commands placed between bodies restore the article\textquotesingle s own printed numbers exactly: Lemma 2.2 and Lemma 2.3 for the two context lemmas, equation (4) for the discrete-time system, equation (8) for the multiplicatively perturbed system, Theorem 5.1 for the counted statement and equations (14), (15), (16), (17) and (18) for the numbered displays of the counted statement and proof. Each of these numbers is the number the version PDF prints for the same object. The two unnumbered displays X4 and X5 remain unnumbered, as in the article.\par
\noindent\textit{Citation map.} The delivered unit invokes exactly three citations: \textbackslash{}cite\{Popova2003\} in the optional argument of each of the two context lemmas (source lines 296--320 region), \textbackslash{}cite[p.\,:239, Proposition 2]\{Lancaster\} and \textbackslash{}cite\{HJ\} inside the counted proof. Their three entries are transcribed from the article\textquotesingle s own in-source thebibliography (source0.tex lines 1821--2017) and print with the article\textquotesingle s own numbers [19], [17] and [14], which are the numbers the version PDF prints for those entries (the article hides six further entries inside \textbackslash{}if0 blocks, so the visible numbering is not the file order). The delivered bibliography contains those three entries and no other, which is exactly the set of keys the delivered text cites.\par
\noindent\textit{Assets.} The e-print of this version is a single gzipped author .tex file; it ships no class, style, bibliography, font or artwork. No retained span contains any \textbackslash{}includegraphics, \textbackslash{}input or \textbackslash{}include command, and the delivered TeX is one self-contained file. Every mathematical symbol, sentence, equation and label of the delivered spans is the authors\textquotesingle{} own native text; no mathematics was written, repaired, re-derived or normalized, and every printed slip inside these spans - including the paper\textquotesingle s own typography such as the run-in \textbackslash{}it statement style and the \textbackslash{}hfill $\square$ proof terminator - is kept literal.\par
\setcounter{equation}{3}
\begin{equation}
x(n+1)=A(n)x(n),\quad n\in \mathbb{N},~x\in \mathbb{R}^{s},  \label{2}
\end{equation}%

\setcounter{section}{2}
\setcounter{theorem}{1}
\begin{lemma}[\protect\cite{Popova2003}]
\label{L3}
\it
Let $V$ be a linear subspace of $\mathbb{R}^{s}$, $\dim V=s-1$;
and let $p\in $ $\mathbb{R}^{s}\backslash V$ be an arbitrary nonzero vector.
If a matrix $H\in \mathbb{R}^{s\times s}$ satisfies the conditions $Hp=p$
and $Hx=0$ for each $x\in V$, then %its norm can be computed by the formula
\begin{equation*}
\| H\| =\dfrac{1}{\sin \sphericalangle (p;V)}.
\end{equation*}
\end{lemma}

\begin{lemma}[\protect\cite{Popova2003}]
\label{L2}
\it
Let $a\colon \mathbb{N}\to \mathbb{R}$ and $b\colon 
\mathbb{N}\to \mathbb{R}$ be arbitrary bounded mappings, and let 
\begin{equation*}
\psi (\mu )=\limsup_{k\to \infty }\bigl(a(k)+\mu b(k)\bigr).
\end{equation*}%
Then the following assertions are valid:

$(1)$ The function $\psi\colon\mathbb{R}\to\mathbb{R}$ is convex and
satisfies the Lipschitz condition on $\mathbb{R}$.

$(2)$ If there exists a strictly increasing sequence $( k_{j})
_{j\in \mathbb{N}}$ of positive integers such that%
\begin{equation*}
\lim_{j\to \infty }a( k_{j}) =\psi ( 0) ,\
\rho \doteq \lim_{j\to \infty }b(k_{j})>0,
\end{equation*}%
then the function $\psi (\cdot )$ is (strictly) monotone increasing on the
interval $[ 0,\infty ) $ and the estimate%
\begin{equation*}
\psi ( \mu ) -\psi ( 0) \ge \rho \mu
\end{equation*}%
is valid for all $\mu \ge 0$. Moreover for each $t\ge 0,$ there exists a $%
\mu _{t},$ $0\le \mu _{t}\le \rho ^{-1}t,$ such that 
\begin{equation*}
\psi ( \mu _{t}) =\psi ( 0) +t.
\end{equation*}
\end{lemma}

\setcounter{equation}{7}
\begin{equation}
x(n+1)=A(n)R(n)x(n),\quad n\in \mathbb{N},\ x\in \mathbb{R}^{s}.  \label{M4}
\end{equation}%
Here the sequence $R(\cdot )$ is called \emph{a multiplicative perturbation%
} whereas system~\eqref{M4} is called  \emph{a multiplicatively perturbed
system}. 

\setcounter{section}{5}
\setcounter{theorem}{0}
\setcounter{equation}{13}
\begin{theorem}
\it
\label{T6}Suppose that system~\eqref{2} has a splitted FSS 
$x_{1}(\cdot),\dots,x_{s}(\cdot)$. Then there exist $\beta>0$ and $\delta >0$ such
that for any $\xi_{i}\in [-\delta,\delta]$, $i=1,\dots ,s$,  there exists an admissible
multiplicative perturbation $R(\cdot )$, satisfying the estimate%
\begin{equation}
\| R-I\| _{\infty }\le \beta \max 
\bigl\{|\xi _{i}| \colon i=1,\dots ,s\bigr\}  
\label{X7}
\end{equation}%
and such that the solutions $\overline{x}_{i}(\cdot )$, $i=1,\ldots ,s$, of
system~\eqref{M4} with the initial conditions $\overline{x}_{i}(1)=x_{i}(1),$
$i=1,\dots ,s$, satisfy the relations%
\begin{equation*}
\lambda [ \overline{x}_{i}] =\lambda [ x_{i}]+\xi _{i},\quad i=1,\dots ,s.
\end{equation*}%
\end{theorem}

\emph{Proof.}
Fix $\sigma=1$. Since the solutions $x_{i}(\cdot )$, 
$i=1,\dots ,s$, are broken away, it follows that there exist a number 
$\gamma \in \left( 0,\frac{\pi }{2}\right] $ and realizing sequences 
$\bigl(k_{j}(i)\bigr)_{j\in \mathbb{N}}\subset \mathbb{N}$
for solutions $x_{i}(\cdot )$, $i=1,\dots ,s$, such that%
\begin{equation*}
\rho _{i}=\lim_{j\to \infty }g_{i}^{\gamma }( k_{j}(i)) >0
\end{equation*}%
for any $i=1,\dots ,s$. Note that the inequality $\rho _{i}\le 1$ is always
valid, since%
\begin{equation*}
\sup \left\{ g^{\gamma}_{i}(k)\colon k\in \mathbb{N},\ i=1,\dots ,s\right\} \le 1.
\end{equation*}%
This, together with Lemma~\ref{L2}, implies that each function 
\begin{equation*}
\Lambda _{i}^{\gamma }( \mu ) \doteq \limsup_{k\to \infty
}( f_{i}(k)+\mu g_{i}^{\gamma }(k))  
\end{equation*}%
satisfies the estimate%
\begin{equation}
\lambda [ x_{i}] +\mu \ge \Lambda _{i}^{\gamma }( \mu
) \ge \lambda [ x_{i}] +\rho \mu ,
\label{X2}
\end{equation}%
where%
\begin{equation*}
\rho =\min \left\{ \rho _{i}\colon i=1,\dots ,s\right\} ,
\end{equation*}%
and for each $t\ge 0$, there exists a $\mu _{t}^{i}\in [ 0,\rho ^{-1}t%
] $ such that%
\begin{equation}
\Lambda _{i}^{\gamma }( \mu _{t}^{i}) =\lambda [ x_{i}]
+t.  \label{X3}
\end{equation}%
Since $\gamma $ is a 
fixed number throughout the proof, from now on we omit the
superscript~$\gamma $.

Let us fix $r\in ( 0,1) $ and $\delta _{1}\in ( 0,\ln (
L_{1}+1) ) $, where%
\begin{equation*}
L_{1}=\frac{r\sin \gamma }{s}
\end{equation*}%
and denote%
\begin{equation*}
L=\frac{L_{1}}{\delta _{1}}.
\end{equation*}%
It is easy to verify that 
\begin{equation*}
\left| \exp ( \delta _{1}) -1\right| <L_{1},\quad
\left| \exp ( -\delta _{1}) -1\right| <L_{1}
\end{equation*}%
and 
\begin{equation*}
\left| \exp ( \tau ) -1\right| \le L\left| \tau
\right| 
\end{equation*}%
for all $\tau \in ( -\infty ,\delta _{1}] $. 
We set 
\begin{equation*}
\delta =\frac{\delta _{1}\rho }{3}
\end{equation*}%
and take an arbitrary $\xi_{i}\in [-\delta,\delta]$, $i=1,\dots,s$. 
Let $\eta\doteq\min\bigl\{\xi_i\colon i=1,\dots,s\bigr\}$ and $\zeta_i\doteq\xi_i-\eta$. Then
$|\eta|\le\delta$ and
$0=\xi_i-\xi_i\le\xi_i-\eta=\zeta_i\le\xi_i+|\eta|\le 2\delta$,
that is,
\begin{equation*}
0\le\zeta_i\le 2\delta.
\end{equation*}
Let
\begin{equation*}
\varepsilon \doteq\max \bigl\{|\xi_{i}| \colon i=1,\dots ,s\bigr\}.
\end{equation*}%
Then 
\begin{equation*}
|\eta| \le \varepsilon  
\label{X4}
\end{equation*}%
and%
\begin{equation*}
\zeta_i\le|\xi_i|+|\eta|\le 2\varepsilon,\quad i=1,\dots,s.   
\label{X5}
\end{equation*}%
Let the quantities $\mu _{i},$ $i=1,\dots ,s,$ be
obtained from the conditions 
\begin{equation*}
\Lambda _{i}(\mu _{i})=\lambda [ x_{i}] +\zeta _{i}.
\end{equation*}%
Since $\zeta_{i}\ge 0$, it follows that the numbers $\mu _{i}$ are well
defined by~\eqref{X3}, and by~\eqref{X2} we have the estimates 
\begin{equation*}
\mu _{i}\ge \Lambda _{i}(\mu _{i})-\lambda [ x_{i}] =\zeta_{i}\ge \rho \mu _{i}\ge 0
\end{equation*}%
for all $i=1,\dots ,s$.


For each $n\in \mathbb{N}$, we introduce a matrix $R(n)\in \mathbb{R}%
^{s\times s}$ by the formulas%
\begin{equation}
R(n)x_{i}( n) =x_{i}( n) \exp \bigl(s_{i}(
n) \bigr),\quad i=1,\dots ,s,  
\label{N7}
\end{equation}%
where 
\begin{equation*}
s_{i}( n) =
\begin{cases}
\eta +\mu _{i} & \textrm{for\ } n\in \Gamma _{i},\\
\eta & \textrm{for\ } n\not\in \Gamma _{i}.
\end{cases}
\end{equation*}%
Then we have the estimates%
\begin{equation*}
|s_{i}(n)|\le|\eta|+\mu_{i}\le|\eta| +\frac{\zeta_{i}}{\rho}\le
\delta+\frac{2\delta}{\rho}=\frac{\delta}{\rho}(\rho+2)\le\frac{3\delta}{\rho}=
\delta_{1}
\end{equation*}%
for $i=1,\dots ,s$ and $n\in \mathbb{N}$. This, together with the definition
of $\delta_{1}$, implies that 
\begin{equation*}
\bigl|\exp\bigl(s_{i}(n)\bigr)-1\bigr| \le L_{1}
\end{equation*}%
and 
\begin{equation*}
\bigl|\exp\bigl(s_{i}(n)\bigr)-1\bigr|\le
L\bigl| s_{i}( n) \bigr| \le 
L\left(|\eta|+\frac{\zeta_i}{\rho}\right)\le
L\left(|\eta|+\frac{2\varepsilon}{\rho}\right)\le
\left( 1+\frac{2}{\rho }%
\right) L\varepsilon.
\end{equation*}%
By the definition of 
a FSS, the vectors $x_{1}(n),\dots,x_{s}(n)$ are linearly independent for each $n\in \mathbb{N}$
and, by~\eqref{N7}, they are eigenvectors of the matrix $R(n)$. This implies
that $R(n)$ is a matrix of simple structure \cite[p.\,239, Proposition 2]{Lancaster} 
and therefore it can be represented by the sum%
\begin{equation*}
R(n)=\sum\limits_{i=1}^{s}P_{n}^{i}\exp\bigl(s_{i}(n)\bigr),
\end{equation*}%
where the root projections $P_{n}^{i}$ are given by the conditions%
\begin{equation*}
P_{n}^{i}x_{i}( n) =x_{i}( n)
\end{equation*}%
and%
\begin{equation*}
P_{n}^{i}x_{j}( n) =0
\end{equation*}%
for $j\neq i$. Moreover 
\begin{equation*}
\sum\limits_{i=1}^{s}P_{n}^{i}=I
\end{equation*}%
for all $n\in \mathbb{N}$. By Lemma~\ref{L3}, we have the estimate 
\begin{equation*}
\left\| P_{n}^{i}\right\| \le \frac{1}{\sin \gamma }
\end{equation*}%
and hence the inequalities 
\begin{equation}
\begin{split}
\left\| R(n)-I\right\| &=  \Bigl\|
\sum\limits_{i=1}^{s}P_{n}^{i}\bigl(\exp\bigl(s_{i}(n)\bigr)-1\bigr)\Bigr\|\\ 
&\le \sum\limits_{i=1}^{s}\left\|
P_{n}^{i}\right\| \left| \exp \bigl( s_{i}( n) \bigr)-1\right| <s\frac{L_{1}}{\sin \gamma }=r. 
\end{split}
\label{X6}
\end{equation}%
Moreover,%
\begin{equation*}
\left\| R(n)-I\right\| \le s\frac{\left| \exp ( s_i(n)) -1\right| }{\sin \gamma }\le \frac{s}{\sin \gamma }\left( 1+%
\frac{2}{\rho }\right) L\varepsilon .
\end{equation*}%
Hence, for the sequence $R(\cdot )=\bigl( R(n)\bigr)_{n\in\mathbb{N}}$ we
have%
\begin{equation*}
\left\| R-I\right\| _{\infty }=\sup_{n\in \mathbb{N}}\left\|
R(n)-I\right\| \le \beta \varepsilon ,
\end{equation*}%
where%
\begin{equation*}
\beta =\frac{Ls\left( 1+\frac{2}{\rho }\right) }{\sin \gamma }.
\end{equation*}%
This proves~\eqref{X7}.


Moreover, we have%
\begin{equation*}
X_{AR}( n+1,n) =X_{A}( n+1,n) R(n)
\end{equation*}%
for all $n\in \mathbb{N}$. Since $r\in ( 0,1) $, then the
condition $\left\| H-I\right\| <r$ implies that $H$ is invertible and%
\begin{equation*}
\left\| H\right\| \le r+1,
\end{equation*}%
\begin{equation*}
\left\| H^{-1}\right\| \le \frac{1}{1-r}
\end{equation*}%
whatever $H\in \mathbb{R}^{s\times s}$ is given~\cite[p.\,301]{HJ}. 
By~\eqref{X6} this
in turn implies that the sequence $R(\cdot )$ is an admissible
multiplicative perturbation.

Consider the FSS $\overline{x}_{i}(\cdot )$, $i=1,\ldots ,s$, of~\eqref{M4}
with such a perturbation with the initial conditions%
\begin{equation*}
\overline{x}_{i}( 1) =x_{i}( 1) ,\quad i=1,\ldots ,s.
\end{equation*}%
For every natural $i\le s$ and $k\ge 2$ we have the equalities
\begin{equation*}
\begin{split}
\overline{x}_{i}(k)&=X_{AR}(k,1)\overline{x}_{i}(1)=X_{AR}(k,1)x_{i}(1)=\prod_{j=1}^{k-1}X_{AR}(j+1,j)x_{i}(1)\\
&=\prod_{j=1}^{k-1}X_{A}(j+1,j)R(j)x_{i}(1)=
\prod_{j=1}^{k-1}X_{A}(j+1,j)\exp\bigl(s_i(j)\bigr)x_{i}(1)\\
&=\exp\bigl(s_i(1)+\ldots+s_{i}(k-1)\bigr)X_A(k,1)x_i(1)=\exp\Bigl(\sum_{j=1}^{k-1}s_i(j)\Bigr)x_i(k).
\end{split}
\end{equation*}
It follows that
the Lyapunov exponents of these solutions satisfy the relations
\begin{equation*}
\begin{split}
\lambda[\overline{x}_{i}] 
&=\limsup_{k\to \infty }\frac{1}{k}\ln\|\overline{x}_{i}(k)\| 
=\limsup_{k\to \infty }\frac{1}{k}\Bigl(\ln\|x_{i}(k)\| +\sum\limits_{j=1}^{k-1}s_{i}(j)\Bigr)\\
&=\limsup_{k\to\infty}\Bigl(f_{i}(k)+\frac{\mu_{i}N_{i}(k-1)}{k}+\frac{\eta(k-1)}{k}\Bigr)\\ 
&=\limsup_{k\to \infty }\Bigl(f_{i}(k)+\frac{\mu_{i}N_{i}(k)}{k}+\frac{\mu_i(N_i(k-1)-N_i(k))}{k}+\eta-\frac{\eta}{k}\Bigr)\\
&=\limsup_{k\to \infty }\bigl(f_{i}(k)+\mu_{i}g_{i}(k)\bigr) +\eta= 
\Lambda _{i}( \mu _{i}) +\eta =\lambda [ x_{i}] +\zeta_{i}+\eta
=\lambda [ x_{i}] +\xi_{i}
\end{split}
\end{equation*}%
for $i=1,\dots ,s$. 
\hfill $\square$

\begin{thebibliography}{99}
\bibitem[14]{HJ} \newblock R.~A. Horn and C.~R. Johnson, \newblock \emph{Matrix Analysis}, \newblock Cambridge University Press, Cambridge, 1985.
\bibitem[17]{Lancaster} \newblock P.~Lancaster and M.~Tismenetsky, \newblock \emph{The Theory of Matrices: with Applications}, \newblock Elsevier, 1985.
\bibitem[19]{Popova2003} \newblock E.~K. Makarov and S.~N. Popova, \newblock Sufficient conditions for the local proportional controllability of {Lyapunov} exponents of linear systems, \newblock \emph{Differ. Equ.}, \textbf{39} (2003), 234--245. [DOI: 10.1023/A:1025157000314]
\end{thebibliography}
\end{document}
